Quan s'il·lumina una superfície metàl·lica amb una radiació ultraviolada $\lambda = 300\ \mathrm{nm}$, el metall emet electrons amb una energia cinètica tan gran que, per a frenar-los (anul·lar el corrent), cal aplicar-hi un potencial de frenada d'1,04~V.

\begin{apartat}{1}
Calculeu l'energia dels fotons incidents i el treball d'extracció (o funció de treball) d'aquest metall.
\end{apartat}

\begin{apartat}{1}
A partir del balanç d'energia de l'efecte fotoelèctric, trobeu l'expressió de la velocitat màxima dels fotoelectrons emesos en funció de la massa dels electrons ($m$), la constant de Planck, la velocitat de la llum, la longitud d'ona de la llum incident i el treball d'extracció ($W_\mathrm{e}$).
\end{apartat}

\dades{%
  $1\ \mathrm{eV} = 1,602\times 10^{-19}\ \mathrm{J}$.\\
  $c = 3,00\times 10^{8}\ \mathrm{m\,s^{-1}}$.\\
  $h = 6,63\times 10^{-34}\ \mathrm{J\,s}$.}
